problem:  
**Conjecture (Natural Reductivity Rigidity for Nonnegatively Curved G.O. Manifolds):**  
Let \((M,g) = (G/H, g)\) be a compact, simply connected geodesic orbit (G.O.) Riemannian manifold with nonnegative sectional curvature. Then \((M,g)\) is naturally reductive. Equivalently, every compact homogeneous G.O. metric with \(\sec \ge 0\) arises from an \(\operatorname{Ad}(H)\)-invariant inner product \(\langle\cdot,\cdot\rangle\) on a reductive complement \(\mathfrak{g} = \mathfrak{h} \oplus \mathfrak{m}\) satisfying  
\[
\langle [X,Y]_{\mathfrak{m}}, Z\rangle + \langle Y, [X,Z]_{\mathfrak{m}}\rangle = 0 \quad \forall\, X,Y,Z \in \mathfrak{m}.
\]

**Background motivation:**  
Known classifications show that:  
- Every **symmetric** space and every **naturally reductive** space is a G.O. manifold.  
- There exist **non‑naturally reductive G.O. metrics** (e.g. certain nilmanifolds or solvmanifolds by Nikonorov), but all known compact examples either have **indefinite curvature** or **mixed sign** sectional curvature.  
- For compact G.O. spaces with \(\sec > 0\), only naturally reductive examples are known (e.g. normal homogeneous spaces or Berger spheres).  

No general theorem currently rules out non‑naturally reductive G.O. metrics with \(\sec \ge 0\); it remains an open structural question whether nonnegative curvature enforces the natural reductivity condition globally.

**Open problem:**  
Prove or disprove the conjecture above. If false, construct an explicit compact, simply connected, non‑naturally reductive G.O. manifold with nonnegative sectional curvature and describe its Lie algebraic structure. If true, characterize geometrically why nonnegative curvature blocks all possible non‑naturally reductive G.O. configurations via algebraic curvature constraints in Nomizu’s formula.

---

Why is it a "good" problem:  
This conjecture targets the precise **boundary between symmetry and curvature** in homogeneous differential geometry. It builds directly on known classification efforts (Kowalski, Vanhecke, Nikonorov) but isolates a point of maximal theoretical tension:

1. **Profound structural insight:** It asks whether the G.O. condition combined with nonnegative curvature automatically enforces the stringent algebraic symmetry characterized by natural reductivity. Establishing this would clarify how curvature restricts the Lie algebraic degrees of freedom in homogeneous metrics.

2. **Bridge across subfields:** The problem connects homogeneous Riemannian geometry, Lie theory (reductive decompositions, isotropy representations), and global curvature conditions. Its resolution likely requires synthesizing algebraic curvature computations (Nomizu-type formulas) with global geometric classification techniques.

3. **Extreme simplicity:** The conjecture can be stated succinctly and understood by anyone familiar with homogeneous spaces. Yet, proving or disproving it would involve deep Lie-algebraic analysis and geometric ingenuity.

4. **Potential impact:**  
   - A positive resolution would unify two important geometric classes (G.O. and naturally reductive manifolds) under curvature constraints, simplifying classification schemes.  
   - A counterexample would expose new phenomena in homogeneous curvature geometry, suggesting undiscovered families of non‑naturally reductive but nonnegatively curved spaces.

Thus, the problem meets all hallmarks of a “good” research problem: **simple statement, profound implications, cross‑disciplinary reach, and location at an uncharted frontier** of differential geometry.